\chapter{實驗結果與討論}\label{chap:results}
\begingroup
\setlength{\emergencystretch}{2em}

\section{統計分析資料與比較範圍}

本章先呈現實驗一的 Baseline 與 Old／New 集成結果，再比較實驗二的特徵選取結果。Baseline 與固定年份集成的驗證方式不同，因此分開呈現；統計檢定只比較資料切分與驗證規則相同的方法。集成與 FS 使用 1999--2011 年訓練 Old 模型，2012--2013 年訓練 New 模型，2014 年分為 DSEL 與 Calibration，2015--2018 年作為 Test。Test 含 12,282 筆公司年度紀錄，其中 287 筆為正類；若按公司計算，共 3,700 家，其中 119 家為破產公司。

固定年份集成與 FS 共使用十個隨機種子及九種特徵設定，合計 90 組設定。統計分析前，先核對各組測試樣本與標籤，再以保存的預測機率及門檻重算指標，確認與原紀錄一致。本次分析沿用原有模型與驗證決策，沒有重新選擇設定。

表中的「合併四年結果（pooled）」先在每個種子內合併四個測試年度計算，再取十個種子的平均值。標準差用來觀察十次執行的分數有多接近。信賴區間與統計檢定則依第~\ref{sec:method-statistics}~節，以公司為單位重抽樣：同一公司的各年度紀錄一起保留，兩個方法使用同一批抽到的公司。全部 165 項比較再一併採 Holm 校正，降低同時比較多種方法時的誤判機會。

本章表格的粗體只幫助辨認指定比較範圍內較高的分數，範圍在各表下方說明。粗體不代表差異已達統計顯著；是否顯著仍依校正後的檢定結果判斷。第三章超參數表的粗體則表示驗證資料選出的參數值，兩者用途不同。

Baseline 的 Old／New 使用 15 組年份邊界，每組包含四種採樣；Retrain 的每種採樣只有一筆完整期間結果。其表格不加上跨種子的標準差，也不把重疊邊界當成獨立重複。與固定年份集成尚未完成同協定的重新調參比較，因此兩類結果的分數差不作為集成效益之檢定證據。

\section{實驗一：Baseline 分類器比較}\label{sec:results-baseline}

表~\ref{tab:results-baseline} 彙整七組已保存的分類器／版本紀錄。Old 與 New 先在同一邊界內平均四種採樣，再平均 15 組邊界；Retrain 直接平均四種採樣。這個彙整方式讓每個邊界與每種採樣具有相同權重，沒有依 Test 挑出最佳組合。各採樣的完整結果列於附錄~\ref{sec:results-baseline-full}。

\begingroup
\small\setstretch{1.12}\setlength{\tabcolsep}{3pt}
\setlength{\LTleft}{\fill}\setlength{\LTright}{\fill}
\begin{longtable}{@{}>{\raggedright\arraybackslash}p{0.34\linewidth}>{\raggedright\arraybackslash}p{0.13\linewidth}>{\raggedright\arraybackslash}p{0.15\linewidth}>{\raggedright\arraybackslash}p{0.15\linewidth}>{\raggedright\arraybackslash}p{0.15\linewidth}@{}}
\caption{實驗一 Baseline 各分類器與訓練策略之測試結果}\label{tab:results-baseline}\\
\toprule
分類器／版本 & 策略 & AUC & F1-score & Recall \\
\midrule\endfirsthead
\multicolumn{5}{c}{表~\thetable（續）}\\
\toprule
分類器／版本 & 策略 & AUC & F1-score & Recall \\
\midrule\endhead
\midrule\multicolumn{5}{r}{續下頁}\\\endfoot
\bottomrule\endlastfoot
XGBoost（較早版本） & Old & 0.7018 & 0.1136 & 0.3365 \\
XGBoost（較早版本） & New & \textbf{0.8526} & \textbf{0.2408} & 0.4689 \\
XGBoost（較早版本） & Retrain & 0.8325 & 0.2117 & \textbf{0.4861} \\
XGBoost（擴充版本） & Old & 0.7091 & 0.1138 & 0.3516 \\
XGBoost（擴充版本） & New & \textbf{0.8589} & \textbf{0.2475} & \textbf{0.4713} \\
XGBoost（擴充版本） & Retrain & 0.8386 & 0.2024 & 0.4686 \\
LightGBM & Old & 0.7212 & 0.1208 & 0.3386 \\
LightGBM & New & \textbf{0.8645} & \textbf{0.2616} & \textbf{0.4762} \\
LightGBM & Retrain & 0.8431 & 0.2206 & 0.4443 \\
Random Forest & Old & 0.7284 & 0.1052 & 0.4258 \\
Random Forest & New & \textbf{0.8453} & \textbf{0.2071} & \textbf{0.4668} \\
Random Forest & Retrain & 0.8294 & 0.1886 & 0.4512 \\
Logistic Regression & Old & 0.6657 & 0.0691 & \textbf{0.6305} \\
Logistic Regression & New & \textbf{0.7074} & \textbf{0.1493} & 0.3710 \\
Logistic Regression & Retrain & 0.6992 & 0.0962 & 0.4817 \\
TabM & Old & 0.6601 & 0.0846 & 0.3800 \\
TabM & New & \textbf{0.8008} & \textbf{0.1860} & \textbf{0.4231} \\
TabM & Retrain & 0.7940 & 0.1496 & 0.4111 \\
TabICL 標記版本 & Old & 0.7062 & 0.0931 & 0.4643 \\
TabICL 標記版本 & New & \textbf{0.8221} & \textbf{0.2057} & 0.4400 \\
TabICL 標記版本 & Retrain & 0.7707 & 0.1380 & \textbf{0.4922} \\
\end{longtable}
\noindent{\footnotesize 資料來源：本研究實驗結果整理。\par Old／New 各先平均同一邊界的四種採樣，再平均 15 組邊界；Retrain 只平均四種採樣。粗體為同一分類器版本內該指標的最高值，四位小數相同者並列；不跨分類器排名，也不表示差異達統計顯著。不同歷史版本的調參與處理條件未完全一致，與固定年份集成須分開解讀。TabICL 為結果標記，後端限制見附錄。}
\endgroup


七組紀錄的 New 平均 AUC 與 F1-score 均高於對應 Old。例如 XGBoost 擴充版本的 Old／New AUC 為 0.7091／0.8589，F1-score 為 0.1138／0.2475；其 Retrain AUC 為 0.8386。這些數值反映本資料中近期期間模型的描述性表現，並未排除樣本量、驗證來源或歷史執行設定差異的影響。

Recall 的方向並不完全相同。Logistic Regression 的 Old Recall 為 0.6305，高於 New 的 0.3710，但 Old F1-score 僅為 0.0691，低於 New 的 0.1493。因此，較高 Recall 不能單獨解釋為較好的分類表現，也需要同時檢視正類預測的精確程度。

XGBoost 兩版本與 TabICL 均依原紀錄分列。歷史調參條件不同，因此僅作描述性比較；後續集成採第三章固定的 XGBoost。

\begin{figure}[htbp]
\centering
\includegraphics[width=\linewidth]{figures/results-baseline-strategies.png}
\caption{Baseline 不同訓練策略之測試結果}\label{fig:results-baseline}
\par\smallskip{\footnotesize 資料來源：本研究依表~\ref{tab:results-baseline}之完整精度結果繪製。Old／New 為十五組邊界與四種採樣之等權平均，Retrain 為四種採樣平均。各指標刻度不同；點位不含信賴區間，分類器依固定順序排列。\par}
\end{figure}

圖~\ref{fig:results-baseline}將相同版本的三種訓練策略放在一起，便於觀察 AUC、F1-score 與 Recall 是否同方向變化。New 在前兩項指標的描述性優勢較一致，Recall 則隨分類器而異。這也說明選擇資料期間時，需要先界定重視的是分數排序、正類涵蓋率，還是精確率與召回率的平衡。

\section{實驗一：Old／New 集成比較}\label{sec:results-experiment1}

\subsection{六個基礎模型與靜態成員數}

本部分使用完整 18 個特徵與固定年份配置。表~\ref{tab:results-pool-models} 先列六個池內模型的測試結果，使後續整合結果可連同成員本身的表現解讀。

\begingroup
\small\setstretch{1.12}\setlength{\tabcolsep}{3pt}
\setlength{\LTleft}{\fill}\setlength{\LTright}{\fill}
\begin{longtable}{@{}>{\raggedright\arraybackslash}p{0.30\linewidth}>{\raggedright\arraybackslash}p{0.21\linewidth}>{\raggedright\arraybackslash}p{0.21\linewidth}>{\raggedright\arraybackslash}p{0.21\linewidth}@{}}
\caption{實驗一六個基礎模型之測試結果}\label{tab:results-pool-models}\\
\toprule
基礎模型 & AUC & F1-score & Recall \\
\midrule\endfirsthead
\multicolumn{4}{c}{表~\thetable（續）}\\
\toprule
基礎模型 & AUC & F1-score & Recall \\
\midrule\endhead
\midrule\multicolumn{4}{r}{續下頁}\\\endfoot
\bottomrule\endlastfoot
Old\_under & 0.7597 & 0.1497 & \textbf{0.3380} \\
Old\_over & 0.6729 & 0.0925 & 0.2882 \\
Old\_hybrid & 0.6803 & 0.0983 & 0.2659 \\
New\_under & \textbf{0.8067} & \textbf{0.2332} & 0.3206 \\
New\_over & 0.7734 & 0.1737 & 0.2557 \\
New\_hybrid & 0.7684 & 0.1724 & 0.2101 \\
\end{longtable}
\noindent{\footnotesize 資料來源：本研究實驗結果整理。\par 完整 18 特徵、固定年份與相同十個種子；under、over、hybrid 分別為 TomekLinks、ADASYN、SMOTEENN。粗體為六模型中該指標最高值，不表示差異達統計顯著。各模型門檻由 Calibration 決定；本表與歷史 Baseline 分開解讀。}
\endgroup


例如 New-under 的 AUC、F1-score 與 Recall 為 0.8067、0.2332 與 0.3206，顯示六個基礎模型的表現並不相同。本表僅列出分數，未另外檢定個別模型與 V-static 或動態方法的差異，因此不由分數高低宣稱差異已達統計顯著。

\begingroup
\small\setstretch{1.12}\setlength{\tabcolsep}{3pt}
\setlength{\LTleft}{\fill}\setlength{\LTright}{\fill}
\begin{longtable}{@{}>{\raggedright\arraybackslash}p{0.15\linewidth}>{\raggedright\arraybackslash}p{0.15\linewidth}>{\raggedright\arraybackslash}p{0.21\linewidth}>{\raggedright\arraybackslash}p{0.21\linewidth}>{\raggedright\arraybackslash}p{0.21\linewidth}@{}}
\caption{實驗一不同靜態成員數之候選平均結果}\label{tab:results-static-size}\\
\toprule
模型數 & 候選數 & AUC & F1-score & Recall \\
\midrule\endfirsthead
\multicolumn{5}{c}{表~\thetable（續）}\\
\toprule
模型數 & 候選數 & AUC & F1-score & Recall \\
\midrule\endhead
\midrule\multicolumn{5}{r}{續下頁}\\\endfoot
\bottomrule\endlastfoot
2 & 9 & 0.7701 & 0.1763 & 0.2348 \\
3 & 18 & 0.7731 & 0.1745 & 0.2468 \\
4 & 15 & 0.7760 & 0.1752 & 0.2440 \\
5 & 6 & 0.7792 & 0.1763 & 0.2469 \\
6 & 1 & \textbf{0.7815} & \textbf{0.1815} & \textbf{0.2690} \\
\end{longtable}
\noindent{\footnotesize 資料來源：本研究實驗結果整理。\par 每個種子先平均同一成員數全部候選的測試分數，再取十個種子平均。粗體為各成員數中該指標最高值，僅供比較，不依此選成員數，也不表示顯著差異。本表沒有把候選再次組成新模型。}
\endgroup


表~\ref{tab:results-static-size} 將 49 個混合候選依成員數分組。兩模型候選的平均 AUC 為 0.7701，六模型為 0.7815；F1-score 並非在每次增加成員時都上升。這些平均值涵蓋不同模型組合，不能解釋為只改變成員數所造成的效果。實際採用的 V-static 仍由 Calibration AP 決定，不依本表的 Test 平均值選擇成員數。

\subsection{集成方法與配對檢定}

先比較使用完整 18 個特徵的集成結果。Old3 與 New3 各整合同一期間的三種採樣模型，Equal6 使用全部六模型；V-static 則依各種子的 Calibration AP，從 49 個混合子集中選定成員。動態方法使用相同六模型池與 DSEL，結果見表~\ref{tab:stats-exp1-descriptive}。

\begingroup
\small\setstretch{1.12}\setlength{\tabcolsep}{3pt}
\setlength{\LTleft}{\fill}\setlength{\LTright}{\fill}
\begin{longtable}{@{}>{\raggedright\arraybackslash}p{0.22\linewidth}>{\raggedright\arraybackslash}p{0.24\linewidth}>{\raggedright\arraybackslash}p{0.24\linewidth}>{\raggedright\arraybackslash}p{0.24\linewidth}@{}}
\caption{實驗一固定年份集成之 pooled 結果}\label{tab:stats-exp1-descriptive}\\
\toprule
方法 & AUC & F1-score & Recall \\
\midrule\endfirsthead
\multicolumn{4}{c}{表~\thetable（續）}\\
\toprule
方法 & AUC & F1-score & Recall \\
\midrule\endhead
\midrule\multicolumn{4}{r}{續下頁}\\\endfoot
\bottomrule\endlastfoot
Old3 & 0.7134\(\pm\)0.0046 & 0.1168\(\pm\)0.0046 & 0.2383\(\pm\)0.0759 \\
New3 & 0.7959\(\pm\)0.0027 & 0.2007\(\pm\)0.0090 & 0.2268\(\pm\)0.0361 \\
Equal6 & 0.7815\(\pm\)0.0026 & 0.1815\(\pm\)0.0065 & 0.2690\(\pm\)0.0414 \\
V-static & \textbf{0.8152}\(\pm\)0.0040 & \textbf{0.2355}\(\pm\)0.0159 & 0.2481\(\pm\)0.0233 \\
DCS-OLA & 0.7386\(\pm\)0.0078 & 0.1542\(\pm\)0.0078 & 0.2909\(\pm\)0.0607 \\
DCS-LCA & 0.7369\(\pm\)0.0087 & 0.1631\(\pm\)0.0083 & \textbf{0.3195}\(\pm\)0.0354 \\
KNORA-E & 0.7613\(\pm\)0.0052 & 0.1837\(\pm\)0.0058 & 0.2714\(\pm\)0.0345 \\
KNORA-U & 0.7481\(\pm\)0.0079 & 0.1817\(\pm\)0.0081 & 0.2620\(\pm\)0.0428 \\
\end{longtable}
\noindent{\footnotesize 資料來源：本研究實驗結果整理。\par 數值為十個種子的平均值與標準差；標準差說明十次執行的分數差異。粗體為同一指標（同一欄）的最高平均值，四位小數相同者並列；不代表差異達統計顯著。V-static 由 Calibration AP 選成員，Equal6 使用全部六模型等權平均。}
\endgroup


\begin{figure}[htbp]
\centering
\includegraphics[width=\linewidth]{figures/results-ensemble-comparison.png}
\caption{未進行 FS 之集成方法比較}\label{fig:results-ensemble}
\par\smallskip{\footnotesize 資料來源：本研究固定年份實驗結果。點為合併四年測試指標的十種子平均，線段為種子標準差，並非公司重抽樣的信賴區間。各指標刻度不同，方法順序不依測試分數排列。\par}
\end{figure}

圖~\ref{fig:results-ensemble}顯示三項指標的相對位置不同。V-static 的平均 AUC 與 F1-score 較高，DCS-LCA 的平均 Recall 較高；線段則顯示各方法對執行隨機性的敏感程度。圖中線段是否重疊不作為顯著性判斷，配對檢定仍依下表及附錄結果。

表~\ref{tab:stats-exp1-auc} 將各方法與 Equal6 比較：差值為正表示 AUC 較高，為負表示較低。95\% 信賴區間表示公司重抽樣估計的差值範圍；是否達到統計顯著，則看 Holm 校正後的 \(p_H\) 是否小於 0.05。

\begingroup
\small\setstretch{1.12}\setlength{\tabcolsep}{3pt}
\setlength{\LTleft}{\fill}\setlength{\LTright}{\fill}
\begin{longtable}{@{}>{\raggedright\arraybackslash}p{0.25\linewidth}>{\raggedright\arraybackslash}p{0.17\linewidth}>{\raggedright\arraybackslash}p{0.33\linewidth}>{\raggedright\arraybackslash}p{0.17\linewidth}@{}}
\caption{實驗一各集成方法相對 Equal6 之 AUC 配對檢定}\label{tab:stats-exp1-auc}\\
\toprule
方法 & \(\Delta\mathrm{AUC}\) & 95\% 信賴區間 & \(p_H\) \\
\midrule\endfirsthead
\multicolumn{4}{c}{表~\thetable（續）}\\
\toprule
方法 & \(\Delta\mathrm{AUC}\) & 95\% 信賴區間 & \(p_H\) \\
\midrule\endhead
\midrule\multicolumn{4}{r}{續下頁}\\\endfoot
\bottomrule\endlastfoot
Old3 & -0.0681 & [-0.0890, -0.0488] & 0.0033\(^{*}\) \\
New3 & +0.0144 & [-0.0019, +0.0313] & 1.0000 \\
V-static & +0.0337 & [+0.0198, +0.0480] & 0.0063\(^{*}\) \\
DCS-OLA & -0.0430 & [-0.0634, -0.0233] & 0.0186\(^{*}\) \\
DCS-LCA & -0.0446 & [-0.0663, -0.0242] & 0.0186\(^{*}\) \\
KNORA-E & -0.0202 & [-0.0387, -0.0029] & 1.0000 \\
KNORA-U & -0.0334 & [-0.0517, -0.0161] & 0.0626 \\
\end{longtable}
\noindent{\footnotesize 資料來源：本研究實驗結果整理。\par 差值為候選方法減比較基準；信賴區間為公司群聚 bootstrap 的逐項 95\% 區間，非同時區間。\(p_H\) 為兩項實驗全部 165 項比較共同採 Holm 校正之雙尾近似 \(p\) 值；星號表示 \(p_H<0.05\)。所有比較使用相同公司重抽權重，模型與驗證門檻固定。}
\endgroup


V-static 的 AUC 與 F1-score 相對 Equal6 分別增加 0.0337 與 0.0540，Holm 校正後 \(p_H\) 分別為 0.00632 與 0.02416。AUC 差值的逐項 95\% 信賴區間為 [0.0198, 0.0480]。其靜態成員由 Calibration 決定，因此該差異屬於驗證選模程序與完整六模型等權平均之比較。相對地，Old3 的 AUC 與 F1-score 差值分別為 \(-0.0681\) 與 \(-0.0647\)，兩者 \(p_H\) 均為 0.00330。

DCS-OLA 與 DCS-LCA 的 AUC 相對 Equal6 分別降低 0.0430 與 0.0446，兩者 \(p_H=0.01860\)。New3、KNORA-E 與 KNORA-U 的 AUC 比較未通過整體 Holm 校正；所有方法相對 Equal6 的 Recall 比較亦未達校正後顯著。這些結果顯示，不能預設動態選擇或跨期間整合必然優於所有控制組。

F1-score 與 Recall 的完整檢定列於附錄~\ref{sec:stats-full}。AUC 為主要報告指標，但其改善不自動代表指定門檻下的 Recall 也同時改善；三項指標須分別解讀。

\subsection{不同測試年度的表現}

\begingroup
\small\setstretch{1.12}\setlength{\tabcolsep}{3pt}
\setlength{\LTleft}{\fill}\setlength{\LTright}{\fill}
\begin{longtable}{@{}>{\raggedright\arraybackslash}p{0.28\linewidth}>{\raggedright\arraybackslash}p{0.12\linewidth}>{\raggedright\arraybackslash}p{0.18\linewidth}>{\raggedright\arraybackslash}p{0.18\linewidth}>{\raggedright\arraybackslash}p{0.18\linewidth}@{}}
\caption{實驗一固定年份集成之逐年測試結果}\label{tab:results-ensemble-yearly}\\
\toprule
方法 & 年度 & AUC & F1-score & Recall \\
\midrule\endfirsthead
\multicolumn{5}{c}{表~\thetable（續）}\\
\toprule
方法 & 年度 & AUC & F1-score & Recall \\
\midrule\endhead
\midrule\multicolumn{5}{r}{續下頁}\\\endfoot
\bottomrule\endlastfoot
Old3 & 2015 & 0.6822 & 0.1289 & 0.2099 \\
Old3 & 2016 & 0.6992 & 0.1301 & 0.2542 \\
Old3 & 2017 & 0.7620 & 0.1032 & 0.2421 \\
Old3 & 2018 & 0.7692 & 0.0928 & 0.2833 \\
New3 & 2015 & 0.7834 & 0.2674 & 0.2550 \\
New3 & 2016 & 0.7731 & 0.2020 & 0.2048 \\
New3 & 2017 & 0.8350 & 0.1956 & 0.2421 \\
New3 & 2018 & 0.8325 & 0.0924 & 0.1667 \\
Equal6 & 2015 & 0.7567 & 0.2275 & 0.2838 \\
Equal6 & 2016 & 0.7560 & 0.1935 & 0.2663 \\
Equal6 & 2017 & 0.8464 & 0.1611 & 0.2596 \\
Equal6 & 2018 & 0.8209 & 0.1071 & 0.2444 \\
V-static & 2015 & \textbf{0.8022} & \textbf{0.2949} & 0.2649 \\
V-static & 2016 & \textbf{0.7937} & \textbf{0.2309} & 0.2205 \\
V-static & 2017 & \textbf{0.8496} & \textbf{0.2451} & 0.2772 \\
V-static & 2018 & \textbf{0.8556} & 0.1304 & 0.2139 \\
DCS-OLA & 2015 & 0.7257 & 0.2000 & 0.2928 \\
DCS-OLA & 2016 & 0.7139 & 0.1629 & 0.2880 \\
DCS-OLA & 2017 & 0.7730 & 0.1232 & 0.2649 \\
DCS-OLA & 2018 & 0.7834 & 0.1089 & 0.3333 \\
DCS-LCA & 2015 & 0.7224 & 0.1909 & 0.2910 \\
DCS-LCA & 2016 & 0.7254 & 0.1826 & \textbf{0.3289} \\
DCS-LCA & 2017 & 0.7567 & 0.1310 & 0.2982 \\
DCS-LCA & 2018 & 0.7790 & \textbf{0.1321} & \textbf{0.4194} \\
KNORA-E & 2015 & 0.7639 & 0.2499 & \textbf{0.2973} \\
KNORA-E & 2016 & 0.7103 & 0.1567 & 0.2133 \\
KNORA-E & 2017 & 0.8016 & 0.1806 & \textbf{0.3088} \\
KNORA-E & 2018 & 0.8111 & 0.1171 & 0.2667 \\
KNORA-U & 2015 & 0.7368 & 0.2355 & 0.2802 \\
KNORA-U & 2016 & 0.7317 & 0.1868 & 0.2482 \\
KNORA-U & 2017 & 0.7806 & 0.1663 & 0.2684 \\
KNORA-U & 2018 & 0.7719 & 0.1009 & 0.2278 \\
\end{longtable}
\noindent{\footnotesize 資料來源：本研究實驗結果整理。\par 各值為十個種子的年度指標平均。粗體表示同一年度八種方法中該指標最高值，不比較不同年度，也不表示差異達統計顯著。2015--2018 年正類紀錄數依序為 111、83、57、36；各年沿用 Calibration 固定的門檻。}
\endgroup


表~\ref{tab:results-ensemble-yearly} 分開呈現四個測試年度。以 V-static 為例，2015 年與 2018 年 AUC 為 0.8022 與 0.8556，F1-score 卻由 0.2949 降至 0.1304。AUC 衡量正負類分數能否排開，F1-score 則與實際分類門檻有關，兩者可能朝不同方向變化。因此，除了合併四年的 AUC，也需要查看逐年結果。

測試正類紀錄數由 2015 年的 111 筆降至 2018 年的 36 筆，類別比例與樣本組成也隨年度不同。年度結果維持原驗證門檻，沒有依測試年度重新校正；本表不將四個年度當成四次獨立實驗，也未據此宣稱變化由特定景氣或漂移因素造成。

\section{實驗二：FS 集成比較}\label{sec:results-experiment2}

接著將八種 FS 條件，與相同種子、測試公司及集成方法的 no-FS 結果配對。兩種保留設定分別使用 9 與 14 個特徵，精簡率為 50.00\% 與 22.22\%。各條件都重新訓練模型，並依相同的 Calibration 規則選擇靜態成員與門檻。

\subsection{特徵條件的完整測試結果}

表~\ref{tab:results-fs-auc}、表~\ref{tab:results-fs-f1} 與表~\ref{tab:results-fs-recall} 分別列出八個 FS 條件及 no-FS 對照。每一欄使用同一種集成規則，每一條件以十個種子彙整，不把不同方法的分數混成單一平均。

\begingroup
\small\setstretch{1.12}\setlength{\tabcolsep}{3pt}
\setlength{\LTleft}{\fill}\setlength{\LTright}{\fill}
\begin{longtable}{@{}>{\raggedright\arraybackslash}p{0.16\linewidth}>{\raggedright\arraybackslash}p{0.123\linewidth}>{\raggedright\arraybackslash}p{0.123\linewidth}>{\raggedright\arraybackslash}p{0.123\linewidth}>{\raggedright\arraybackslash}p{0.123\linewidth}>{\raggedright\arraybackslash}p{0.123\linewidth}>{\raggedright\arraybackslash}p{0.123\linewidth}@{}}
\caption{實驗二各特徵條件之 pooled AUC}\label{tab:results-fs-auc}\\
\toprule
FS & Equal6 & V-static & OLA & LCA & KNORA-E & KNORA-U \\
\midrule\endfirsthead
\multicolumn{7}{c}{表~\thetable（續）}\\
\toprule
FS & Equal6 & V-static & OLA & LCA & KNORA-E & KNORA-U \\
\midrule\endhead
\midrule\multicolumn{7}{r}{續下頁}\\\endfoot
\bottomrule\endlastfoot
no-FS（18） & \shortstack[l]{0.7815\\{\scriptsize(\(\pm\)0.0026)}} & \shortstack[l]{\textbf{0.8152}\\{\scriptsize(\(\pm\)0.0040)}} & \shortstack[l]{0.7386\\{\scriptsize(\(\pm\)0.0078)}} & \shortstack[l]{0.7369\\{\scriptsize(\(\pm\)0.0087)}} & \shortstack[l]{\textbf{0.7613}\\{\scriptsize(\(\pm\)0.0052)}} & \shortstack[l]{0.7481\\{\scriptsize(\(\pm\)0.0079)}} \\
MI 50\% & \shortstack[l]{0.7644\\{\scriptsize(\(\pm\)0.0059)}} & \shortstack[l]{0.7826\\{\scriptsize(\(\pm\)0.0119)}} & \shortstack[l]{0.7129\\{\scriptsize(\(\pm\)0.0146)}} & \shortstack[l]{0.6938\\{\scriptsize(\(\pm\)0.0219)}} & \shortstack[l]{0.7339\\{\scriptsize(\(\pm\)0.0105)}} & \shortstack[l]{0.7262\\{\scriptsize(\(\pm\)0.0067)}} \\
MI 80\% & \shortstack[l]{0.7822\\{\scriptsize(\(\pm\)0.0038)}} & \shortstack[l]{0.8145\\{\scriptsize(\(\pm\)0.0039)}} & \shortstack[l]{0.7319\\{\scriptsize(\(\pm\)0.0059)}} & \shortstack[l]{0.7185\\{\scriptsize(\(\pm\)0.0054)}} & \shortstack[l]{0.7514\\{\scriptsize(\(\pm\)0.0075)}} & \shortstack[l]{0.7558\\{\scriptsize(\(\pm\)0.0092)}} \\
CART 50\% & \shortstack[l]{0.7511\\{\scriptsize(\(\pm\)0.0137)}} & \shortstack[l]{0.7931\\{\scriptsize(\(\pm\)0.0101)}} & \shortstack[l]{0.7087\\{\scriptsize(\(\pm\)0.0230)}} & \shortstack[l]{0.7074\\{\scriptsize(\(\pm\)0.0238)}} & \shortstack[l]{0.7178\\{\scriptsize(\(\pm\)0.0168)}} & \shortstack[l]{0.7252\\{\scriptsize(\(\pm\)0.0120)}} \\
CART 80\% & \shortstack[l]{0.7812\\{\scriptsize(\(\pm\)0.0056)}} & \shortstack[l]{0.8140\\{\scriptsize(\(\pm\)0.0032)}} & \shortstack[l]{0.7412\\{\scriptsize(\(\pm\)0.0089)}} & \shortstack[l]{\textbf{0.7372}\\{\scriptsize(\(\pm\)0.0094)}} & \shortstack[l]{0.7557\\{\scriptsize(\(\pm\)0.0083)}} & \shortstack[l]{0.7494\\{\scriptsize(\(\pm\)0.0071)}} \\
SHAP 50\% & \shortstack[l]{0.7595\\{\scriptsize(\(\pm\)0.0024)}} & \shortstack[l]{0.7971\\{\scriptsize(\(\pm\)0.0022)}} & \shortstack[l]{0.7353\\{\scriptsize(\(\pm\)0.0086)}} & \shortstack[l]{0.7261\\{\scriptsize(\(\pm\)0.0072)}} & \shortstack[l]{0.7398\\{\scriptsize(\(\pm\)0.0073)}} & \shortstack[l]{0.7341\\{\scriptsize(\(\pm\)0.0080)}} \\
SHAP 80\% & \shortstack[l]{0.7815\\{\scriptsize(\(\pm\)0.0027)}} & \shortstack[l]{0.8117\\{\scriptsize(\(\pm\)0.0040)}} & \shortstack[l]{0.7344\\{\scriptsize(\(\pm\)0.0069)}} & \shortstack[l]{0.7279\\{\scriptsize(\(\pm\)0.0047)}} & \shortstack[l]{0.7484\\{\scriptsize(\(\pm\)0.0078)}} & \shortstack[l]{0.7468\\{\scriptsize(\(\pm\)0.0043)}} \\
RFE 50\% & \shortstack[l]{0.7617\\{\scriptsize(\(\pm\)0.0087)}} & \shortstack[l]{0.8011\\{\scriptsize(\(\pm\)0.0105)}} & \shortstack[l]{0.7107\\{\scriptsize(\(\pm\)0.0109)}} & \shortstack[l]{0.7083\\{\scriptsize(\(\pm\)0.0125)}} & \shortstack[l]{0.7235\\{\scriptsize(\(\pm\)0.0091)}} & \shortstack[l]{0.7356\\{\scriptsize(\(\pm\)0.0126)}} \\
RFE 80\% & \shortstack[l]{\textbf{0.7837}\\{\scriptsize(\(\pm\)0.0069)}} & \shortstack[l]{0.8150\\{\scriptsize(\(\pm\)0.0067)}} & \shortstack[l]{\textbf{0.7500}\\{\scriptsize(\(\pm\)0.0115)}} & \shortstack[l]{0.7356\\{\scriptsize(\(\pm\)0.0076)}} & \shortstack[l]{0.7600\\{\scriptsize(\(\pm\)0.0115)}} & \shortstack[l]{\textbf{0.7568}\\{\scriptsize(\(\pm\)0.0104)}} \\
\end{longtable}
\noindent{\footnotesize 資料來源：本研究實驗結果整理。\par 各格上列為十個種子的平均值，下列為種子標準差。粗體表示同一方法（同一欄）九種特徵條件中最高的平均分數，四位小數相同者並列；不代表已選定最佳設定或差異達統計顯著。50\%、80\% 分別保留 9、14 個特徵，no-FS 為 18 個；V-static 可依各條件的 Calibration 重新選成員。}
\endgroup

\begingroup
\small\setstretch{1.12}\setlength{\tabcolsep}{3pt}
\setlength{\LTleft}{\fill}\setlength{\LTright}{\fill}
\begin{longtable}{@{}>{\raggedright\arraybackslash}p{0.16\linewidth}>{\raggedright\arraybackslash}p{0.123\linewidth}>{\raggedright\arraybackslash}p{0.123\linewidth}>{\raggedright\arraybackslash}p{0.123\linewidth}>{\raggedright\arraybackslash}p{0.123\linewidth}>{\raggedright\arraybackslash}p{0.123\linewidth}>{\raggedright\arraybackslash}p{0.123\linewidth}@{}}
\caption{實驗二各特徵條件之 pooled F1-score}\label{tab:results-fs-f1}\\
\toprule
FS & Equal6 & V-static & OLA & LCA & KNORA-E & KNORA-U \\
\midrule\endfirsthead
\multicolumn{7}{c}{表~\thetable（續）}\\
\toprule
FS & Equal6 & V-static & OLA & LCA & KNORA-E & KNORA-U \\
\midrule\endhead
\midrule\multicolumn{7}{r}{續下頁}\\\endfoot
\bottomrule\endlastfoot
no-FS（18） & \shortstack[l]{0.1815\\{\scriptsize(\(\pm\)0.0065)}} & \shortstack[l]{\textbf{0.2355}\\{\scriptsize(\(\pm\)0.0159)}} & \shortstack[l]{0.1542\\{\scriptsize(\(\pm\)0.0078)}} & \shortstack[l]{\textbf{0.1631}\\{\scriptsize(\(\pm\)0.0083)}} & \shortstack[l]{0.1837\\{\scriptsize(\(\pm\)0.0058)}} & \shortstack[l]{0.1817\\{\scriptsize(\(\pm\)0.0081)}} \\
MI 50\% & \shortstack[l]{0.1647\\{\scriptsize(\(\pm\)0.0093)}} & \shortstack[l]{0.1960\\{\scriptsize(\(\pm\)0.0173)}} & \shortstack[l]{0.1522\\{\scriptsize(\(\pm\)0.0241)}} & \shortstack[l]{0.1354\\{\scriptsize(\(\pm\)0.0174)}} & \shortstack[l]{0.1594\\{\scriptsize(\(\pm\)0.0113)}} & \shortstack[l]{0.1821\\{\scriptsize(\(\pm\)0.0073)}} \\
MI 80\% & \shortstack[l]{0.1818\\{\scriptsize(\(\pm\)0.0121)}} & \shortstack[l]{0.2263\\{\scriptsize(\(\pm\)0.0171)}} & \shortstack[l]{0.1542\\{\scriptsize(\(\pm\)0.0155)}} & \shortstack[l]{0.1449\\{\scriptsize(\(\pm\)0.0161)}} & \shortstack[l]{0.1780\\{\scriptsize(\(\pm\)0.0073)}} & \shortstack[l]{0.1922\\{\scriptsize(\(\pm\)0.0141)}} \\
CART 50\% & \shortstack[l]{0.1541\\{\scriptsize(\(\pm\)0.0123)}} & \shortstack[l]{0.1942\\{\scriptsize(\(\pm\)0.0187)}} & \shortstack[l]{0.1258\\{\scriptsize(\(\pm\)0.0207)}} & \shortstack[l]{0.1300\\{\scriptsize(\(\pm\)0.0130)}} & \shortstack[l]{0.1441\\{\scriptsize(\(\pm\)0.0109)}} & \shortstack[l]{0.1621\\{\scriptsize(\(\pm\)0.0110)}} \\
CART 80\% & \shortstack[l]{0.1765\\{\scriptsize(\(\pm\)0.0061)}} & \shortstack[l]{0.2321\\{\scriptsize(\(\pm\)0.0130)}} & \shortstack[l]{0.1465\\{\scriptsize(\(\pm\)0.0090)}} & \shortstack[l]{0.1555\\{\scriptsize(\(\pm\)0.0128)}} & \shortstack[l]{\textbf{0.1857}\\{\scriptsize(\(\pm\)0.0129)}} & \shortstack[l]{0.1909\\{\scriptsize(\(\pm\)0.0098)}} \\
SHAP 50\% & \shortstack[l]{0.1638\\{\scriptsize(\(\pm\)0.0078)}} & \shortstack[l]{0.1867\\{\scriptsize(\(\pm\)0.0197)}} & \shortstack[l]{0.1500\\{\scriptsize(\(\pm\)0.0114)}} & \shortstack[l]{0.1361\\{\scriptsize(\(\pm\)0.0143)}} & \shortstack[l]{0.1648\\{\scriptsize(\(\pm\)0.0079)}} & \shortstack[l]{0.1708\\{\scriptsize(\(\pm\)0.0088)}} \\
SHAP 80\% & \shortstack[l]{0.1849\\{\scriptsize(\(\pm\)0.0070)}} & \shortstack[l]{0.2203\\{\scriptsize(\(\pm\)0.0204)}} & \shortstack[l]{0.1367\\{\scriptsize(\(\pm\)0.0185)}} & \shortstack[l]{0.1274\\{\scriptsize(\(\pm\)0.0104)}} & \shortstack[l]{0.1718\\{\scriptsize(\(\pm\)0.0098)}} & \shortstack[l]{0.1920\\{\scriptsize(\(\pm\)0.0091)}} \\
RFE 50\% & \shortstack[l]{0.1577\\{\scriptsize(\(\pm\)0.0178)}} & \shortstack[l]{0.2119\\{\scriptsize(\(\pm\)0.0225)}} & \shortstack[l]{0.1372\\{\scriptsize(\(\pm\)0.0120)}} & \shortstack[l]{0.1392\\{\scriptsize(\(\pm\)0.0118)}} & \shortstack[l]{0.1502\\{\scriptsize(\(\pm\)0.0116)}} & \shortstack[l]{0.1697\\{\scriptsize(\(\pm\)0.0225)}} \\
RFE 80\% & \shortstack[l]{\textbf{0.1873}\\{\scriptsize(\(\pm\)0.0142)}} & \shortstack[l]{0.2297\\{\scriptsize(\(\pm\)0.0166)}} & \shortstack[l]{\textbf{0.1604}\\{\scriptsize(\(\pm\)0.0118)}} & \shortstack[l]{0.1556\\{\scriptsize(\(\pm\)0.0142)}} & \shortstack[l]{0.1780\\{\scriptsize(\(\pm\)0.0137)}} & \shortstack[l]{\textbf{0.1945}\\{\scriptsize(\(\pm\)0.0144)}} \\
\end{longtable}
\noindent{\footnotesize 資料來源：本研究實驗結果整理。\par 各格上列為十個種子的平均值，下列為種子標準差。粗體表示同一方法（同一欄）九種特徵條件中最高的平均分數，四位小數相同者並列；不代表已選定最佳設定或差異達統計顯著。50\%、80\% 分別保留 9、14 個特徵，no-FS 為 18 個；V-static 可依各條件的 Calibration 重新選成員。}
\endgroup

\begingroup
\small\setstretch{1.12}\setlength{\tabcolsep}{3pt}
\setlength{\LTleft}{\fill}\setlength{\LTright}{\fill}
\begin{longtable}{@{}>{\raggedright\arraybackslash}p{0.16\linewidth}>{\raggedright\arraybackslash}p{0.123\linewidth}>{\raggedright\arraybackslash}p{0.123\linewidth}>{\raggedright\arraybackslash}p{0.123\linewidth}>{\raggedright\arraybackslash}p{0.123\linewidth}>{\raggedright\arraybackslash}p{0.123\linewidth}>{\raggedright\arraybackslash}p{0.123\linewidth}@{}}
\caption{實驗二各特徵條件之 pooled Recall}\label{tab:results-fs-recall}\\
\toprule
FS & Equal6 & V-static & OLA & LCA & KNORA-E & KNORA-U \\
\midrule\endfirsthead
\multicolumn{7}{c}{表~\thetable（續）}\\
\toprule
FS & Equal6 & V-static & OLA & LCA & KNORA-E & KNORA-U \\
\midrule\endhead
\midrule\multicolumn{7}{r}{續下頁}\\\endfoot
\bottomrule\endlastfoot
no-FS（18） & \shortstack[l]{0.2690\\{\scriptsize(\(\pm\)0.0414)}} & \shortstack[l]{0.2481\\{\scriptsize(\(\pm\)0.0233)}} & \shortstack[l]{0.2909\\{\scriptsize(\(\pm\)0.0607)}} & \shortstack[l]{\textbf{0.3195}\\{\scriptsize(\(\pm\)0.0354)}} & \shortstack[l]{0.2714\\{\scriptsize(\(\pm\)0.0345)}} & \shortstack[l]{0.2620\\{\scriptsize(\(\pm\)0.0428)}} \\
MI 50\% & \shortstack[l]{\textbf{0.3530}\\{\scriptsize(\(\pm\)0.0546)}} & \shortstack[l]{0.2763\\{\scriptsize(\(\pm\)0.0667)}} & \shortstack[l]{0.2902\\{\scriptsize(\(\pm\)0.0914)}} & \shortstack[l]{0.2376\\{\scriptsize(\(\pm\)0.0994)}} & \shortstack[l]{\textbf{0.2895}\\{\scriptsize(\(\pm\)0.0547)}} & \shortstack[l]{\textbf{0.2812}\\{\scriptsize(\(\pm\)0.0405)}} \\
MI 80\% & \shortstack[l]{0.2509\\{\scriptsize(\(\pm\)0.0665)}} & \shortstack[l]{0.3300\\{\scriptsize(\(\pm\)0.0353)}} & \shortstack[l]{0.2533\\{\scriptsize(\(\pm\)0.0866)}} & \shortstack[l]{0.2098\\{\scriptsize(\(\pm\)0.0793)}} & \shortstack[l]{0.2634\\{\scriptsize(\(\pm\)0.0614)}} & \shortstack[l]{0.2679\\{\scriptsize(\(\pm\)0.0577)}} \\
CART 50\% & \shortstack[l]{0.2425\\{\scriptsize(\(\pm\)0.0628)}} & \shortstack[l]{0.2878\\{\scriptsize(\(\pm\)0.0632)}} & \shortstack[l]{0.2613\\{\scriptsize(\(\pm\)0.1378)}} & \shortstack[l]{0.1348\\{\scriptsize(\(\pm\)0.0285)}} & \shortstack[l]{0.2470\\{\scriptsize(\(\pm\)0.0406)}} & \shortstack[l]{0.2101\\{\scriptsize(\(\pm\)0.0609)}} \\
CART 80\% & \shortstack[l]{0.2240\\{\scriptsize(\(\pm\)0.0523)}} & \shortstack[l]{0.2819\\{\scriptsize(\(\pm\)0.0334)}} & \shortstack[l]{0.2115\\{\scriptsize(\(\pm\)0.0790)}} & \shortstack[l]{0.2429\\{\scriptsize(\(\pm\)0.0692)}} & \shortstack[l]{0.2638\\{\scriptsize(\(\pm\)0.0437)}} & \shortstack[l]{0.2578\\{\scriptsize(\(\pm\)0.0596)}} \\
SHAP 50\% & \shortstack[l]{0.2624\\{\scriptsize(\(\pm\)0.0716)}} & \shortstack[l]{\textbf{0.3502}\\{\scriptsize(\(\pm\)0.0602)}} & \shortstack[l]{0.2178\\{\scriptsize(\(\pm\)0.0530)}} & \shortstack[l]{0.2418\\{\scriptsize(\(\pm\)0.1389)}} & \shortstack[l]{0.2401\\{\scriptsize(\(\pm\)0.0462)}} & \shortstack[l]{0.2732\\{\scriptsize(\(\pm\)0.0585)}} \\
SHAP 80\% & \shortstack[l]{0.2672\\{\scriptsize(\(\pm\)0.0427)}} & \shortstack[l]{0.3171\\{\scriptsize(\(\pm\)0.0402)}} & \shortstack[l]{\textbf{0.3540}\\{\scriptsize(\(\pm\)0.0995)}} & \shortstack[l]{0.2279\\{\scriptsize(\(\pm\)0.1395)}} & \shortstack[l]{0.2533\\{\scriptsize(\(\pm\)0.0292)}} & \shortstack[l]{0.2756\\{\scriptsize(\(\pm\)0.0235)}} \\
RFE 50\% & \shortstack[l]{0.2732\\{\scriptsize(\(\pm\)0.0908)}} & \shortstack[l]{0.2805\\{\scriptsize(\(\pm\)0.0626)}} & \shortstack[l]{0.2394\\{\scriptsize(\(\pm\)0.0924)}} & \shortstack[l]{0.2927\\{\scriptsize(\(\pm\)0.0359)}} & \shortstack[l]{0.2631\\{\scriptsize(\(\pm\)0.0621)}} & \shortstack[l]{0.2188\\{\scriptsize(\(\pm\)0.0725)}} \\
RFE 80\% & \shortstack[l]{0.2425\\{\scriptsize(\(\pm\)0.0450)}} & \shortstack[l]{0.2997\\{\scriptsize(\(\pm\)0.0517)}} & \shortstack[l]{0.2655\\{\scriptsize(\(\pm\)0.1069)}} & \shortstack[l]{0.2829\\{\scriptsize(\(\pm\)0.0745)}} & \shortstack[l]{0.2596\\{\scriptsize(\(\pm\)0.0557)}} & \shortstack[l]{0.2564\\{\scriptsize(\(\pm\)0.0412)}} \\
\end{longtable}
\noindent{\footnotesize 資料來源：本研究實驗結果整理。\par 各格上列為十個種子的平均值，下列為種子標準差。粗體表示同一方法（同一欄）九種特徵條件中最高的平均分數，四位小數相同者並列；不代表已選定最佳設定或差異達統計顯著。50\%、80\% 分別保留 9、14 個特徵，no-FS 為 18 個；V-static 可依各條件的 Calibration 重新選成員。}
\endgroup


以 Equal6 為例，no-FS、MI 50\% 與 MI 80\% 的 AUC 為 0.7815、0.7644 與 0.7822，顯示大幅刪減與保留較多特徵的結果不同。V-static 在 SHAP 50\% 下的 Recall 為 0.3502，高於 no-FS 的 0.2481，但 F1-score 為 0.1867，低於 no-FS 的 0.2355。因而不能把少數類涵蓋率的增加直接視為所有指標均改善。

\subsection{與同方法 no-FS 的配對比較}

\begingroup
\small\setstretch{1.12}\setlength{\tabcolsep}{3pt}
\setlength{\LTleft}{\fill}\setlength{\LTright}{\fill}
\begin{longtable}{@{}>{\raggedright\arraybackslash}p{0.16\linewidth}>{\raggedright\arraybackslash}p{0.123\linewidth}>{\raggedright\arraybackslash}p{0.123\linewidth}>{\raggedright\arraybackslash}p{0.123\linewidth}>{\raggedright\arraybackslash}p{0.123\linewidth}>{\raggedright\arraybackslash}p{0.123\linewidth}>{\raggedright\arraybackslash}p{0.123\linewidth}@{}}
\caption{實驗二各 FS 設定相對同方法 no-FS 之 AUC 差值}\label{tab:stats-exp2-auc}\\
\toprule
FS & Equal6 & V-static & OLA & LCA & KNORA-E & KNORA-U \\
\midrule\endfirsthead
\multicolumn{7}{c}{表~\thetable（續）}\\
\toprule
FS & Equal6 & V-static & OLA & LCA & KNORA-E & KNORA-U \\
\midrule\endhead
\midrule\multicolumn{7}{r}{續下頁}\\\endfoot
\bottomrule\endlastfoot
MI 50\% & -0.0171 & -0.0327 & -0.0257 & -0.0431 & -0.0274 & -0.0220 \\
MI 80\% & +0.0007 & -0.0007 & -0.0066 & -0.0184 & -0.0100 & +0.0076 \\
CART 50\% & -0.0304 & -0.0221 & -0.0299 & -0.0295 & -0.0435\(^{*}\) & -0.0229 \\
CART 80\% & -0.0003 & -0.0013 & +0.0026 & +0.0003 & -0.0056 & +0.0013 \\
SHAP 50\% & -0.0220 & -0.0182 & -0.0033 & -0.0109 & -0.0216 & -0.0141 \\
SHAP 80\% & -0.0000 & -0.0036 & -0.0041 & -0.0090 & -0.0129 & -0.0013 \\
RFE 50\% & -0.0198 & -0.0141 & -0.0278 & -0.0286 & -0.0378 & -0.0125 \\
RFE 80\% & +0.0022 & -0.0002 & +0.0115 & -0.0013 & -0.0014 & +0.0087 \\
\end{longtable}
\noindent{\footnotesize 資料來源：本研究實驗結果整理。\par 50\% 與 80\% 設定實際保留 9 與 14 個特徵。正值表示 AUC 增加。星號定義與表~\ref{tab:stats-exp1-auc} 相同；完整區間、F1-score、Recall 與校正值見附錄~\ref{sec:stats-full}。}
\endgroup


\begin{figure}[htbp]
\centering
\includegraphics[width=\linewidth]{figures/results-fs-auc-differences.png}
\caption{各 FS 條件相對同方法 no-FS 之 AUC 差值}\label{fig:results-fs-differences}
\par\smallskip{\footnotesize 資料來源：本研究固定年份配對結果。每格為十種子平均的 FS 減 no-FS，正值表示增加、負值表示下降。色彩只表示差值方向與大小，不表示統計顯著；50\%／80\% 分別實際保留 9／14 個特徵。\par}
\end{figure}

圖~\ref{fig:results-fs-differences}保留全部 48 組 AUC 差值。相同 FS 條件在不同集成方法下可能有不同方向，保留較多特徵也不一定每格都優於保留較少特徵。此圖用來閱讀變化位置，是否有足夠證據支持改善，仍須看共同校正後的檢定。

表~\ref{tab:stats-exp2-auc} 呈現 FS 與同方法 no-FS 的 AUC 差值。完整的 48 組 AUC 比較，以及 F1-score、Recall 的檢定結果與信賴區間，列於附錄~\ref{sec:stats-full}。分數增加不一定達到統計顯著；沒有達到顯著，也不能反過來認定兩者表現相同。

Holm 校正後，沒有 FS 條件呈現顯著的 AUC 增加，F1-score 的全部 48 項比較也未達顯著。CART 50\% 搭配 KNORA-E 的 AUC 低 0.0435，\(p_H=0.03900\)，達到本次設定的顯著門檻。不過，這個值靠近 0.05，且隨機重抽估計仍有波動，因此對這一項結果保留較審慎的解讀。

Recall 共有九項比較通過 Holm 校正，其中 V-static 在 MI 80\%、SHAP 50\%、SHAP 80\% 及 RFE 80\% 下增加，其餘五項為 DCS 方法下的下降。例如，SHAP 50\% 的 V-static Recall 增加 0.1021，\(p_H=0.00632\)；CART 50\% 的 DCS-LCA Recall 降低 0.1847，\(p_H=0.00330\)。此差異顯示 FS 的影響與集成方法有關，不能僅因部分 Recall 改善，就推論 FS 全面提升區辨能力或精確率與召回率的平衡。

V-static 在 FS 後會重新選擇成員與門檻，因此分數變化包含整個選模流程的影響。其他方法雖然保持集成規則，仍會重新訓練與決定門檻。故本章比較的是加入 FS 後整套流程的結果，無法把所有變化都歸因於特徵數減少。

\section{研究問題的綜合討論}\label{sec:results-discussion}

\subsection{訓練資料期間與單一分類器}

對研究問題一，七組歷史 Baseline 紀錄均呈現 New 平均 AUC 與 F1-score 高於 Old，顯示資料期間值得列為建模條件。但這項結果仍同時包含訓練筆數、驗證資料與歷史參數的差異，無法單獨估計資料新近程度的因果效果。Retrain 提供合併資料的描述性基準，並未在固定年份集成協定下重新建立可直接檢定的對照。

因此，本章可以說明近期資料模型在這批紀錄中的表現，卻不據此主張所有應用都應刪除歷史資料，也不將 Baseline 表中最高的測試分數作為後續選型依據。XGBoost 的共同基底定位維持第三章所述的固定研究設定。

\subsection{集成規則與模型成員}

對研究問題二，V-static 相對 Equal6 的 AUC 與 F1-score 增加，在本次公司配對分析與共同校正下仍有統計證據。這表示於本資料中，由 Calibration 選擇混合子集的完整程序，比一律保留全部六模型更合適。由於 V-static 本身要求同時含 Old 與 New，且本研究未對它與 New3 作同一組正式配對檢定，不能進一步說成所有跨期間集成都優於近期資料集成。

動態方法亦未一致超越 Equal6。DCS 的局部選擇會受到 DSEL 與特徵距離影響，但目前沒有單獨操弄 DSEL 大小或近鄰數的實驗，無法由結果斷定是哪一環節造成差異。實務上的啟示是保留簡單的等權與單一期間控制組，讓動態方法的價值由同條件結果判斷，而非由方法名稱或複雜度推論。

\subsection{特徵精簡與少數類辨識}

對研究問題三，FS 將輸入由 18 個特徵減至 9 或 14 個，但共同校正後並未出現顯著的 AUC 增加，F1-score 也沒有顯著改善。部分條件的 Recall 增加，另有條件下降，顯示 FS 的影響需要連同集成規則及門檻解讀。沒有顯著下降的條件，也不能直接宣稱已達到與完整特徵相同的效果；若要證明可接受的效能損失，仍需另定等效或非劣性界線。

本研究記錄了特徵數與入選次數，但未以同一測量程序比較所有條件的訓練及推論時間，因此只能確認輸入維度減少，不能換算成已證實的時間或成本節省。特徵選取的價值也不宜只用最高分數衡量，而應連同精簡程度、正類漏判與誤報的取捨一起評估。

\section{統計推論範圍與限制}

本章的統計分析是在原實驗完成後補做，用來檢查已保存的分數差異。全部 165 項比較均保留，不只列出分數增加或達到顯著者。重抽樣一起保留同公司的各年度紀錄，但沒有重新訓練模型，因此不能代表換一批訓練資料或換到其他年份時的結果。不同公司也可能受同一產業或景氣影響，這部分未另行估計。

本章使用公開資料中的公司層級狀態，該標籤在同一公司所有年度保持相同，且公司可能跨擬合與測試期間出現。因此，結論限於本資料與既定時間分區下的回溯性公司狀態分類。統計差異不等同模型已具備提前一年預測能力，也不能據此宣稱對全新公司或不同市場具有普遍優越性。第三章的欄位表已依實際發布版修正，但模型使用的財務欄位、數值與預測結果保持原狀。
\endgroup
